\documentclass[11pt,oneside,openany]{book}

% =============================================================================
% Configuración general
% =============================================================================
\usepackage[utf8]{inputenc}
\usepackage[T1]{fontenc}
\usepackage[provide=*,spanish]{babel}
\usepackage{lmodern}
\usepackage[a4paper,top=2.30cm,bottom=2.45cm,left=2.45cm,right=2.20cm,
            headheight=15pt,footskip=1.25cm]{geometry}
\linespread{1.06}
\setlength{\parskip}{0.55em}
\setlength{\parindent}{0pt}
\setlength{\emergencystretch}{2em}

% =============================================================================
% Matemática, tablas, figuras y recursos visuales
% =============================================================================
\usepackage{amsmath,amssymb}
\usepackage{graphicx}
\graphicspath{{assets/}}
\usepackage{xcolor}
\usepackage{booktabs}
\usepackage{array}
\usepackage{tabularx}
\usepackage{longtable}
\usepackage{float}
\usepackage{hyperref}

\definecolor{AzulProfundo}{HTML}{102A43}
\definecolor{AzulMedio}{HTML}{176B87}
\definecolor{Turquesa}{HTML}{2A9D8F}
\definecolor{TurquesaOscuro}{HTML}{207A70}
\definecolor{Arena}{HTML}{F4E9D8}
\definecolor{ArenaClara}{HTML}{FBF6EE}
\definecolor{CelesteClaro}{HTML}{EAF4F8}
\definecolor{VerdeClaro}{HTML}{EAF7F3}
\definecolor{AmarilloClaro}{HTML}{FFF6DF}
\definecolor{GrisTexto}{HTML}{3F4D5A}
\definecolor{GrisClaro}{HTML}{EEF2F5}
\definecolor{RojoSuave}{HTML}{C95F5F}

\hypersetup{
  colorlinks=true,
  linkcolor=AzulMedio,
  citecolor=TurquesaOscuro,
  urlcolor=AzulMedio,
  pdftitle={Revisión bibliográfica y selección de algoritmos para MAP},
  pdfauthor={Jorge Gabriel Palacios Sales; Pablo Daniel Barillas Moreno; Roberto Emiliano Otoniel Camposeco Torres; Abby Sofia Donis Agreda},
  pdfsubject={Modelos de aprendizaje automático para identificar errores conceptuales matemáticos}
}

\makeatletter
\def\@makechapterhead#1{%
  \vspace*{18pt}%
  {\parindent\z@\raggedright\normalfont
    \ifnum\c@secnumdepth>\m@ne
      \sffamily\bfseries\color{Turquesa}\fontsize{36}{40}\selectfont
      \thechapter\par\nobreak\vskip 8pt
    \fi
    \sffamily\bfseries\color{AzulProfundo}\Huge #1\par\nobreak
    \vskip 7pt\color{Turquesa}\hrule height 1.1pt\vskip 28pt}}
\def\@makeschapterhead#1{%
  \vspace*{18pt}%
  {\parindent\z@\raggedright\normalfont
    \sffamily\bfseries\color{AzulProfundo}\Huge #1\par\nobreak
    \vskip 7pt\color{Turquesa}\hrule height 1.1pt\vskip 28pt}}
\renewcommand\section{\@startsection{section}{1}{\z@}%
  {-3.5ex \@plus -1ex \@minus -.2ex}{2.3ex \@plus .2ex}%
  {\normalfont\Large\sffamily\bfseries\color{AzulMedio}}}
\renewcommand\subsection{\@startsection{subsection}{2}{\z@}%
  {-3.25ex\@plus -1ex \@minus -.2ex}{1.5ex \@plus .2ex}%
  {\normalfont\large\sffamily\bfseries\color{TurquesaOscuro}}}
\renewcommand\subsubsection{\@startsection{subsubsection}{3}{\z@}%
  {-3.25ex\@plus -1ex \@minus -.2ex}{1.5ex \@plus .2ex}%
  {\normalfont\normalsize\sffamily\bfseries\color{AzulProfundo}}}
\makeatother

% =============================================================================
% Componentes reutilizables
% =============================================================================
\newsavebox{\infobox}
\newenvironment{hallazgo}[1][]{%
  \begin{lrbox}{\infobox}\begin{minipage}{0.92\textwidth}
  \color{GrisTexto}\sffamily\textbf{#1}\par\smallskip\rmfamily}
  {\end{minipage}\end{lrbox}\par\medskip\noindent
   \colorbox{CelesteClaro}{\usebox{\infobox}}\par\medskip}
\newenvironment{decision}[1][]{%
  \begin{lrbox}{\infobox}\begin{minipage}{0.92\textwidth}
  \color{GrisTexto}\sffamily\textbf{#1}\par\smallskip\rmfamily}
  {\end{minipage}\end{lrbox}\par\medskip\noindent
   \colorbox{VerdeClaro}{\usebox{\infobox}}\par\medskip}
\newenvironment{advertencia}[1][]{%
  \begin{lrbox}{\infobox}\begin{minipage}{0.92\textwidth}
  \color{GrisTexto}\sffamily\textbf{#1}\par\smallskip\rmfamily}
  {\end{minipage}\end{lrbox}\par\medskip\noindent
   \colorbox{AmarilloClaro}{\usebox{\infobox}}\par\medskip}
\newenvironment{ficha}[1][]{%
  \begin{lrbox}{\infobox}\begin{minipage}{0.90\textwidth}
  \color{AzulProfundo}\sffamily\textbf{#1}\par\smallskip\rmfamily}
  {\end{minipage}\end{lrbox}\par\medskip\noindent
   \fcolorbox{Arena}{ArenaClara}{\usebox{\infobox}}\par\medskip}

\newcolumntype{Y}{>{\raggedright\arraybackslash}X}
\newcolumntype{C}[1]{>{\centering\arraybackslash}p{#1}}
\newcolumntype{P}[1]{>{\raggedright\arraybackslash}p{#1}}
\newcommand{\codigo}[1]{\texttt{#1}}
\newcommand{\articulo}[1]{\textcolor{AzulMedio}{\textbf{#1}}}
\newcommand{\nivel}[1]{\textcolor{TurquesaOscuro}{\sffamily\bfseries #1}}
\newcommand{\figura}[4][0.90\textwidth]{%
  \begin{figure}[htbp]
    \centering
    \includegraphics[width=#1]{#2}
    \caption{#3}
    \label{#4}
  \end{figure}}

\begin{document}
\hypersetup{pageanchor=false}

% =============================================================================
% Portada
% =============================================================================
\begin{titlepage}
\pagecolor{AzulProfundo}
\color{white}

\begin{minipage}[t]{0.35\textwidth}
  \includegraphics[width=3.2cm]{logo_uvg.png}
\end{minipage}%
\begin{minipage}[t]{0.63\textwidth}
  \raggedleft
  {\sffamily\bfseries\Large Universidad del Valle de Guatemala\par}
  \vspace{1mm}
  {\sffamily Facultad de Ingeniería\par}
  {\sffamily Departamento de Ciencias de la Computación\par}
\end{minipage}

\vfill

{\sffamily\bfseries\color{Turquesa}\Large Proyecto 2 -- Aprendizaje automático\par}
\vspace{0.55cm}
{\sffamily\bfseries\fontsize{31}{36}\selectfont
Revisión bibliográfica y\par
selección preliminar de modelos\par}
\vspace{0.6cm}
{\Large\color{Arena}
Identificación de errores conceptuales matemáticos en explicaciones
estudiantiles mediante NLP\par}

\vspace{1.0cm}
\noindent\fcolorbox{white!40!AzulProfundo}{AzulProfundo}{%
\begin{minipage}{0.94\textwidth}
\vspace{3mm}
\color{white}
\begin{tabularx}{\textwidth}{@{}>{\sffamily\bfseries}p{3.5cm}Y@{}}
Caso de estudio: & MAP -- Charting Student Math Misunderstandings\\[1.5mm]
Curso: & CC3084 -- Data Science\\[1.5mm]
Docente: & Lynette García Pérez\\[1.5mm]
Grupo y sección: & Grupo 1, sección 10\\[1.5mm]
Integrantes: & Jorge Gabriel Palacios Sales\\
& Pablo Daniel Barillas Moreno\\
& Roberto Emiliano Otoniel Camposeco Torres\\
& Abby Sofia Donis Agreda\\[1.5mm]
Área: & Procesamiento del Lenguaje Natural y analítica educativa\\
\end{tabularx}
\vspace{3mm}
\end{minipage}}

\vfill
{\sffamily\large Guatemala, agosto de 2026\par}
\end{titlepage}
\nopagecolor
\hypersetup{pageanchor=true}

\frontmatter

\chapter*{Resumen}
\addcontentsline{toc}{chapter}{Resumen}

Este documento revisa la literatura sobre calificación automática de respuestas
cortas, clasificación de explicaciones estudiantiles y detección de errores
conceptuales, con énfasis en modelos aplicables al reto \emph{MAP -- Charting
Student Math Misunderstandings}. El corpus local analizado contiene 36,696
explicaciones asociadas con 15 preguntas, 49 opciones distintas, 35 concepciones
erróneas y 65 etiquetas compuestas. Esta combinación hace que el problema no sea una simple
clasificación de texto: la interpretación depende de la pregunta, la respuesta
seleccionada, la corrección matemática y la distribución de clases de cola larga.

La evidencia revisada muestra una evolución desde similitud léxica, ingeniería de
características y máquinas de vectores de soporte, hacia redes recurrentes,
Transformers ajustados y adaptación de dominio. Los hallazgos convergen en cuatro
ideas: (a) los modelos lineales con TF--IDF siguen siendo líneas base competitivas,
baratas e interpretables; (b) incluir la pregunta y el contexto supera el uso
aislado de la respuesta; (c) los Transformers capturan paráfrasis y relaciones
semánticas, pero requieren validación estricta y control de sesgos; y (d) la
adaptación al dominio y los ensambles híbridos suelen mejorar la robustez cuando
el corpus etiquetado es limitado o desbalanceado.

Se selecciona preliminarmente una estrategia de cuatro niveles. Primero, un
baseline de TF--IDF de palabras y caracteres con regresión logística y SVM lineal.
Segundo, un modelo principal basado en DeBERTa/BERT que reciba conjuntamente
pregunta, opción y explicación, con cabezas jerárquicas para \codigo{Category} y
\codigo{Misconception}. Tercero, adaptación de dominio mediante preentrenamiento
continuado o MathBERT. Cuarto, un ensamble calibrado que produzca las tres parejas
\codigo{Category:Misconception} mejor posicionadas, coherente con MAP@3. Los LLM
generativos se reservan como comparación y herramienta de auditoría, no como
solución primaria, por reproducibilidad, costo, privacidad y trazabilidad.

\begin{decision}[Decisión preliminar]
La ruta recomendada es \textbf{baseline lineal híbrido $\rightarrow$ Transformer
contextual y multitarea $\rightarrow$ adaptación de dominio $\rightarrow$
ensamble y ranking top--3}. La evaluación debe comparar una partición estratificada
para preguntas conocidas con otra agrupada por \codigo{QuestionId} para transferencia;
ambas deben reportar MAP@3, macro F1, desempeño por clase/pregunta y calibración.
\end{decision}

\chapter*{Alcance y trazabilidad}
\addcontentsline{toc}{chapter}{Alcance y trazabilidad}

\begin{tabularx}{\textwidth}{@{}P{3.3cm}Y@{}}
\toprule
\textbf{Elemento} & \textbf{Tratamiento en este informe}\\
\midrule
Revisión bibliográfica & Artículos de revistas arbitradas e indexadas como corpus
principal; trabajos fundacionales de conferencia o preprint se identifican como
fuentes complementarias.\\
Norma de citación & Citas autor--fecha en el cuerpo y referencias finales en estilo
APA, séptima edición.\\
Selección preliminar & Algoritmos, arquitectura, variables de entrada, validación,
métricas, riesgos y criterios de descarte.\\
Datos del proyecto & Resultados ya obtenidos en el análisis exploratorio; el
notebook no fue reejecutado para preparar este documento.\\
Límite & No se reportan resultados de entrenamiento. Las puntuaciones de la matriz
de decisión son juicios de diseño, no métricas experimentales.\\
\bottomrule
\end{tabularx}

\tableofcontents
\listoftables
\listoffigures

\mainmatter

% =============================================================================
\chapter{Planteamiento del problema}
% =============================================================================

\section{Contexto educativo y computacional}

Los errores matemáticos no son únicamente respuestas incorrectas. Con frecuencia
constituyen evidencias de una regla aplicada fuera de contexto, una confusión de
signos, una generalización excesiva o una representación incompleta del problema.
La explicación escrita permite diferenciar un descuido de cálculo de una idea
equivocada persistente. Esta distinción es pedagógicamente útil porque una
retroalimentación adecuada depende de la causa del error y no solo del resultado.
El análisis de errores puede convertir una respuesta fallida en evidencia del
razonamiento del estudiante y en una oportunidad de aprendizaje más profundo
(Rushton, 2018).

El reto MAP plantea automatizar parte de ese diagnóstico (Kaggle, 2025). A partir de una
pregunta de opción múltiple, la opción elegida y una explicación breve, un sistema
debe predecir la categoría de la respuesta y, cuando corresponde, el error
conceptual asociado. El resultado esperado no es texto libre, sino una lista
ordenada de hasta tres pares \codigo{Category:Misconception}. Por tanto, el
problema combina clasificación multiclase, predicción jerárquica, ranking y
procesamiento del lenguaje matemático.

\begin{hallazgo}[Situación problemática]
Revisar manualmente miles de explicaciones breves exige tiempo especializado y
produce variabilidad entre evaluadores. Un clasificador confiable puede priorizar
casos, sugerir etiquetas y acelerar la retroalimentación; no debe sustituir de
forma ciega el juicio docente, especialmente en clases raras o explicaciones
ambiguas.
\end{hallazgo}

\section{Problema científico}

\textbf{Pregunta de investigación.} ¿Qué combinación de representaciones léxicas,
modelos contextuales y estrategias de adaptación al dominio permite ordenar con
mayor precisión las categorías y los errores conceptuales matemáticos presentes
en explicaciones estudiantiles cortas, manteniendo generalización a preguntas no
vistas, reproducibilidad y trazabilidad?

La pregunta contiene cinco exigencias. Primero, comprender lenguaje informal y
notación matemática. Segundo, condicionar la explicación a la pregunta y a la
opción seleccionada. Tercero, aprender clases desbalanceadas. Cuarto, producir un
ranking top--3 y no una sola etiqueta. Quinto, demostrar que el modelo reconoce
patrones conceptuales y no memoriza identificadores o frases de una pregunta.

\section{Objetivos}

\subsection{Objetivo general}

Revisar críticamente algoritmos de aprendizaje automático aplicados a calificación
automática, clasificación de respuestas cortas y detección de errores conceptuales,
para seleccionar una arquitectura preliminar adecuada al reto MAP.

\subsection{Objetivos específicos}

\begin{enumerate}
  \item Identificar familias de modelos, representaciones y esquemas de evaluación
  utilizados en literatura académica comparable.
  \item Contrastar sus ventajas y limitaciones frente al tamaño, estructura,
  desbalance y dependencia contextual del corpus MAP.
  \item Seleccionar líneas base, un modelo principal, alternativas de adaptación de
  dominio y un método de ensamble/ranking.
  \item Diseñar un plan experimental que evite fuga entre preguntas y mida
  desempeño global, por clase y por pregunta.
  \item Establecer controles de confiabilidad, explicabilidad y supervisión humana.
\end{enumerate}

% =============================================================================
\chapter{Metodología de revisión bibliográfica}
% =============================================================================

\section{Alcance de la búsqueda}

La revisión es narrativa y orientada al diseño. Se priorizaron publicaciones entre
2015 y agosto de 2026 en revistas arbitradas e indexadas relacionadas con
inteligencia artificial en educación, tecnologías de aprendizaje, aprendizaje
interactivo, aplicaciones de IA y educación científica. Se consultaron las páginas
oficiales de editoriales y los registros DOI para verificar título, año, volumen y
paginación o número de artículo.

Las cadenas conceptuales de búsqueda combinaron términos como \emph{automatic
short-answer grading}, \emph{student response classification}, \emph{math
misconception detection}, \emph{BERT}, \emph{Transformer}, \emph{domain-specific
pretraining}, \emph{question context}, \emph{augmentation}, \emph{robustness} y
\emph{human-in-the-loop}. Se incluyeron trabajos que aportaran al menos uno de
estos elementos: comparación de algoritmos, representación textual, uso de
contexto, adaptación de dominio, transferencia, evaluación de generalización o
análisis de confiabilidad.

\section{Criterios de inclusión y exclusión}

\begin{table}[H]
\centering
\caption{Criterios aplicados en la revisión}
\label{tab:criterios-revision}
\begin{tabularx}{\textwidth}{@{}P{3.2cm}YY@{}}
\toprule
\textbf{Dimensión} & \textbf{Inclusión} & \textbf{Exclusión o uso secundario}\\
\midrule
Tipo de fuente & Artículo completo en revista arbitrada; metadatos verificables y
DOI cuando existe. & Blogs, resúmenes comerciales y repositorios sin revisión como
evidencia principal.\\
Problema & Respuestas cortas, clasificación de razonamiento, retroalimentación o
errores conceptuales. & Calificación de ensayos largos sin conexión metodológica
con respuestas breves.\\
Método & Modelos clásicos, redes neuronales, Transformers, LLM, ensambles,
adaptación y robustez. & Sistemas sin descripción suficiente del método o la
evaluación.\\
Transferibilidad & Resultados que permitan derivar una decisión para MAP. & Cifras
aisladas cuya tarea o métrica no es comparable.\\
Fuentes fundacionales & Conferencias arbitradas y preprints se usan para documentar
BERT, DeBERTa, MathBERT y el origen del benchmark. & No se confunden con el núcleo
de artículos de revista solicitado.\\
\bottomrule
\end{tabularx}
\end{table}

\section{Procedimiento de síntesis}

Para cada trabajo se registraron: tarea, datos, familia algorítmica, incorporación
de contexto, métricas, aporte y limitación transferible. Después se agruparon los
estudios por generación metodológica y se contrastaron con las restricciones del
corpus local. Las cifras publicadas se conservan como evidencia de cada escenario,
pero no se comparan directamente entre artículos porque cambian corpus, idiomas,
escalas de nota y particiones.

\begin{advertencia}[Interpretación responsable]
Una mejora de exactitud en otro conjunto no garantiza la misma mejora en MAP. La
revisión sirve para formular hipótesis y reducir el espacio de búsqueda; solamente
la evaluación local, con particiones sin fuga, puede confirmar la selección.
\end{advertencia}

% =============================================================================
\chapter{El caso MAP y sus requisitos de modelado}
% =============================================================================

\section{Unidad de análisis y variables}

Cada observación representa una explicación estudiantil vinculada con una
pregunta y una opción. La señal útil está distribuida entre variables: el texto
describe el razonamiento, la opción acota el resultado elegido y la pregunta
determina el significado matemático de las palabras y símbolos. Esto sugiere una
entrada compuesta y no un clasificador que lea solo la explicación.

\begin{table}[H]
\centering
\caption{Traducción de las variables de MAP a decisiones de modelado}
\label{tab:variables-modelo}
\begin{tabularx}{\textwidth}{@{}P{3.7cm}YY@{}}
\toprule
\textbf{Señal} & \textbf{Función} & \textbf{Tratamiento propuesto}\\
\midrule
\codigo{QuestionId}/pregunta & Define el contexto conceptual y el vocabulario
relevante. & Incorporar texto o identificador controlado; agrupar por pregunta en
validación.\\
\codigo{AnswerText} u opción & Expresa el resultado que el estudiante defendió. &
Concatenar con la explicación mediante separadores.\\
\codigo{StudentExplanation} & Contiene las pistas lingüísticas, operaciones y
negaciones. & Conservar números, signos, fracciones y negaciones; producir
representación léxica y contextual.\\
\codigo{Category} & Describe la clase general de corrección/atribución. & Primera
cabeza de clasificación o primer nivel jerárquico.\\
\codigo{Misconception} & Etiqueta conceptual específica, válida solo en ciertos
casos. & Segunda cabeza condicional; máscara de pérdida cuando no aplica.\\
Pareja de salida & Unidad evaluada y ordenada. & Probabilidad conjunta y ranking
de tres candidatos.\\
\bottomrule
\end{tabularx}
\end{table}

\section{Hallazgos del análisis exploratorio local}

El análisis previo identificó 36,696 observaciones, 15 preguntas, 49 opciones
distintas, 35 concepciones erróneas y 65 etiquetas compuestas. La categoría
\codigo{True\_Correct} representa 40.34\%. La explicación típica es breve: la
mediana es de 13 palabras y la media de 15.36. Además, 26,836 filas (73.1\%) no contienen una etiqueta en
\codigo{Misconception}; esta ausencia es principalmente estructural, porque una
etiqueta específica no corresponde a todas las categorías, y no debe imputarse
como si fuera un valor perdido ordinario.

\figura[0.88\textwidth]{distribucion_misconceptions.png}{Distribución de etiquetas
de error conceptual en el conjunto analizado. La cola larga exige métricas y
estrategias sensibles a clases minoritarias.}{fig:misconceptions}

La Figura~\ref{fig:misconceptions} muestra que unas pocas etiquetas concentran
muchas observaciones y varias poseen soporte reducido. En estas condiciones, la
exactitud puede ocultar el fracaso en conceptos raros. Se requieren macro F1,
recall por clase y análisis de errores, además de MAP@3.

\figura[0.90\textwidth]{cruces_por_pregunta.png}{Cruces por pregunta obtenidos en
el análisis exploratorio. El patrón varía entre contextos y justifica comparar
validación estratificada y agrupada.}{fig:preguntas}

La heterogeneidad entre preguntas observada en la Figura~\ref{fig:preguntas}
implica riesgo de fuga: una partición aleatoria por filas puede colocar frases muy
similares de una misma pregunta en entrenamiento y validación. Un resultado alto
en esa partición mediría, en parte, memorización del contexto. Por ello se propone
reportar dos escenarios: estratificado por filas para preguntas conocidas y
agrupado por \codigo{QuestionId} como prueba de transferencia. Debido a que 32 de
las 35 concepciones son exclusivas de una pregunta, ninguno debe interpretarse de
forma aislada.

\section{Requisitos derivados}

\begin{enumerate}
  \item \textbf{Contextualidad:} interpretar explicación, pregunta y opción en
  conjunto.
  \item \textbf{Sensibilidad matemática:} conservar símbolos, cantidades,
  operadores, signos y negaciones.
  \item \textbf{Jerarquía:} modelar primero la categoría y después el concepto
  erróneo cuando corresponda.
  \item \textbf{Cola larga:} ponderar clases, usar muestreo controlado y evitar que
  las clases frecuentes dominen la pérdida.
  \item \textbf{Ranking:} calibrar probabilidades y optimizar una lista top--3.
  \item \textbf{Generalización:} probar preguntas no vistas y variaciones
  lingüísticas.
  \item \textbf{Trazabilidad:} disponer de una línea base interpretable y análisis
  humano de falsos positivos/negativos.
\end{enumerate}

% =============================================================================
\chapter{Revisión de algoritmos y arquitecturas}
% =============================================================================

\section{De reglas y similitud a aprendizaje supervisado}

La revisión histórica de Burrows et al. (2015), que organiza más de ochenta
trabajos de calificación automática de respuestas cortas, distingue una evolución
desde reglas y coincidencia con respuestas modelo hacia métodos estadísticos y de
aprendizaje automático. Los primeros sistemas eran transparentes, pero dependían
de plantillas y vocabularios elaborados por expertos. Funcionaban mejor cuando la
respuesta esperada estaba estrechamente delimitada y peor ante paráfrasis,
ortografía no estándar o razonamientos novedosos.

Los métodos de similitud comparan la respuesta estudiantil con una referencia por
superposición de palabras, distancia de edición, similitud coseno o recursos
semánticos. Son útiles como características y como control de calidad, pero no
constituyen por sí solos una solución suficiente para MAP: dos estudiantes pueden
usar vocabulario parecido y defender conclusiones conceptualmente opuestas, o usar
palabras distintas para expresar el mismo error.

\section{TF--IDF, regresión logística y SVM}

La representación TF--IDF asigna mayor peso a términos informativos y reduce la
influencia de palabras frecuentes. Con $tf(t,d)$ como frecuencia del término $t$
en el documento $d$, y $N$ documentos, una forma común es

\[
\operatorname{tfidf}(t,d)=tf(t,d)\log\left(\frac{N+1}{df(t)+1}\right).
\]

Sobre esta matriz dispersa, la regresión logística estima probabilidades de clase
y la SVM lineal maximiza márgenes de separación. Los n-gramas de palabras capturan
frases como ``se cambia el signo'', mientras los n-gramas de caracteres toleran
errores ortográficos y conservan fragmentos de notación. La revisión de GradeAid
señala que SVM, bolsa de palabras y \emph{word embeddings} han sido recursos muy
frecuentes en ASAG; su propuesta híbrida combina TF--IDF y señales de BERT con
regresores y logra errores competitivos en varios conjuntos públicos (del Gobbo
et al., 2023).

Para MAP estos métodos ofrecen tres ventajas: entrenamiento rápido, coeficientes
inspeccionables y un piso reproducible. Su limitación es la semántica: una bolsa de
palabras no representa bien el orden, la referencia contextual o el efecto de una
negación distante. Por ello deben considerarse baseline obligatorio y componente
de ensamble, no única solución final.

\section{Ingeniería de características y ensambles}

Sahu y Bhowmick (2020) combinaron características léxicas, de alineamiento y
semánticas mediante un esquema de regresión apilada. El ensamble mejoró la
calificación automática frente a modelos individuales, lo que sugiere que ninguna
representación captura por sí sola toda la evidencia de una respuesta breve. La
lección transferible es conservar señales complementarias: n-gramas, longitud,
coincidencias numéricas, operadores, negaciones, opción elegida y representación
contextual.

En MAP un ensamble no necesita ser excesivamente complejo. Puede combinar la
probabilidad calibrada del modelo lineal y la del Transformer. Si $p_L(y\mid x)$
es la distribución lineal y $p_T(y\mid x)$ la contextual, una primera formulación
es

\[
p_E(y\mid x)=\alpha p_T(y\mid x)+(1-\alpha)p_L(y\mid x),
\quad 0\leq\alpha\leq1,
\]

donde $\alpha$ se elige exclusivamente con validación. También puede aprenderse un
meta-modelo sobre predicciones \emph{out-of-fold}; entrenarlo sobre predicciones
del mismo conjunto produciría fuga.

\section{Modelos recurrentes y orden de palabras}

Las redes LSTM modelan secuencias y mitigan parte del desvanecimiento del gradiente
de RNN simples. Zhang et al. (2022) propusieron un modelo para preguntas
semiabiertas que integra información general y específica del dominio con LSTM.
Su evaluación, basada en más de 16,000 respuestas de siete preguntas, respalda la
utilidad de conservar orden y contexto. Sin embargo, las LSTM procesan tokens
secuencialmente y suelen depender más de la cantidad de datos etiquetados que un
Transformer preentrenado.

Se mantendría una BiLSTM como ablación si el costo lo permite, pero no es la
primera candidata: el corpus MAP posee textos cortos y suficientes señales para
comparar un baseline lineal fuerte con un Transformer. Esa comparación responde
mejor a la pregunta práctica de cuánto valor aporta la contextualización moderna.

\section{La pregunta y el estudiante como contexto}

Zhang et al. (2020) demostraron que combinar modelos de la pregunta y del
estudiante puede mejorar métodos basados únicamente en la respuesta. Aunque su
arquitectura utiliza una red de creencia profunda, el aporte más relevante para
MAP es conceptual: la misma frase cambia de interpretación según el enunciado y el
historial o contexto disponible.

MAP no proporciona necesariamente una secuencia longitudinal completa por
estudiante, pero sí ofrece \codigo{QuestionId} y opción. En consecuencia, se puede
preservar la idea mediante una entrada estructurada:

\begin{center}
\fcolorbox{Arena}{ArenaClara}{\parbox{0.88\textwidth}{\centering\ttfamily
[PREGUNTA] texto o identificador [SEP] [OPCIÓN] respuesta elegida [SEP]
[EXPLICACIÓN] razonamiento estudiantil}}
\end{center}

El identificador no debe convertirse en un atajo. Cuando el texto completo de la
pregunta esté disponible, es preferible porque ofrece significado transferible.
Si solo existe el ID, debe comprobarse el desempeño en preguntas no vistas y
considerarse una codificación ``desconocida''.

\section{BERT y Transformers ajustados}

BERT aprende representaciones bidireccionales mediante preentrenamiento con texto
no etiquetado y después se ajusta a una tarea supervisada (Devlin et al., 2019).
Zhu et al. (2022) combinaron BERT, BiLSTM, redes cápsula y una función de pérdida
específica, superando métodos previos en conjuntos de referencia de ASAG. El
resultado apoya la capacidad de los Transformers para reconocer equivalencia
semántica con pocos ejemplos en comparación con redes entrenadas desde cero.

La revisión de Morley et al. (2026) sobre los años fundacionales de Transformers
en calificación automática de respuestas científicas encontró que BERT y sus
variantes son la familia predominante y que el contexto de pregunta, rúbrica o
material de apoyo puede mejorar el sistema. Al mismo tiempo, la revisión señala
problemas pendientes de confiabilidad, explicación y sesgo. Por ello no basta con
reportar una métrica agregada.

DeBERTa modifica la atención y la codificación posicional mediante atención
desentrelazada y un decodificador mejorado (He et al., 2021). Se propone como
candidato principal por su fortaleza general en clasificación, pero debe
compararse con BERT/RoBERTa bajo el mismo presupuesto. La elección definitiva no
se realizará por popularidad, sino por MAP@3, macro F1, estabilidad entre semillas
y costo.

\section{Adaptación de dominio para matemáticas}

El lenguaje matemático escolar contiene expresiones y símbolos cuya distribución
difiere del texto general. MathBERT fue preentrenado sobre material de educación
matemática y mostró la viabilidad de especializar un Transformer para tareas del
dominio (Shen et al., 2021). Como se trata de un preprint, se utiliza como fuente
arquitectónica complementaria y no como evidencia principal de revista.

La evidencia indexada de Bonthu et al. (2024) es más generalizable: su marco
PyDSBERT construye un corpus no etiquetado del dominio, continúa el preentrenamiento
y ajusta el modelo a calificación. En su escenario reportó 88.77\% de exactitud y
91.59\% de F1, con mejoras de varios puntos sobre referencias. Estas cifras no se
trasladan directamente a MAP, pero justifican probar \emph{domain-adaptive
pretraining} cuando exista texto matemático adicional y el presupuesto lo permita.

Wang et al. (2025) extendieron BERT con una perspectiva ontológica para clasificar
niveles de experiencia científica. Su combinación de conocimiento de dominio,
balanceo y evaluación mediante F1, exactitud y kappa ponderado confirma que la
estructura conceptual puede complementar la representación contextual. Para MAP,
la taxonomía de errores puede actuar como conocimiento: descripciones de las
etiquetas, relaciones entre conceptos y restricciones de compatibilidad con la
pregunta pueden alimentar un clasificador o un re-ranker.

\section{Transferencia y aumento de datos}

Bonthu et al. (2023) ajustaron modelos de oraciones y aplicaron reemplazo de
sinónimos e intercambio aleatorio; reportaron aumentos de 4.91\% en exactitud y
3.36\% en F1 en su tarea. La transferencia es atractiva para clases raras, pero el
aumento debe respetar la semántica matemática. Cambiar un número, signo, operador
o negación puede alterar la etiqueta correcta. En lugar de aumentos ciegos se
recomiendan transformaciones conservadoras: correcciones ortográficas verificadas,
variantes de puntuación, espacios alrededor de operadores y paráfrasis revisadas
por una persona con conocimiento matemático.

\section{Generalización a preguntas no vistas}

Padó et al. (2024) observaron que los Transformers eran robustos cuando las
respuestas se parecían al dominio visto, pero no superaron necesariamente modelos
lingüísticos diseñados a mano en tareas con preguntas no vistas. Este hallazgo es
central para MAP: un modelo puede parecer preciso en una división aleatoria y
fallar cuando cambia el concepto o el enunciado.

La consecuencia metodológica es directa. La evaluación debe comparar una división
estratificada con \emph{GroupKFold} o una variante de dejar preguntas fuera. Deben
separarse dos objetivos: desempeño operativo sobre preguntas conocidas y capacidad
de transferencia. El segundo escenario contiene numerosas etiquetas no vistas por
construcción; por ello es una prueba de extrapolación y no reemplaza la evaluación
operativa.

\section{LLM generativos y evaluación con intervención humana}

Los LLM pueden clasificar mediante instrucciones y ejemplos sin entrenamiento
convencional. Grévisse (2024) evaluó respuestas médicas breves en 12 cursos y tres
idiomas; GPT--4 y Gemini alcanzaron acuerdos moderados, pero la supervisión humana
siguió siendo necesaria. El benchmark específico de errores matemáticos de Otero
et al. (2025), formado por 55 conceptos erróneos y 220 ejemplos diagnósticos,
encontró que restringir el tema y añadir revisión de educadores puede elevar el
desempeño de modelos grandes hasta 83.9\% en su configuración. La lección no es que
el LLM deba reemplazar al clasificador, sino que la información temática y el
juicio experto reducen ambigüedad.

Para el proyecto, un LLM alojado externamente introduce variabilidad de versión,
costo, latencia, privacidad y dificultad para repetir exactamente el experimento.
Puede usarse como comparación \emph{zero-shot}/\emph{few-shot}, generador de
hipótesis o asistente para revisar errores, siempre sin enviar datos sensibles y
registrando modelo, fecha, parámetros y plantilla de prompt.

\section{Confiabilidad, ataques y supervisión}

Schneider et al. (2023) estudiaron calificación multilingüe y multiformato a gran
escala. Su trabajo muestra que combinar Transformers con intervención humana puede
alcanzar niveles comparables a asistentes de docencia y, sobre todo, permite
controlar tipos de error. Esto respalda un diseño de apoyo a decisiones: aceptar
automáticamente solo predicciones bien calibradas y enviar casos inciertos a
revisión.

Filighera et al. (2024) mostraron que insertar adjetivos y adverbios en respuestas
incorrectas puede reducir entre 10 y 22 puntos la exactitud de sistemas BERT/T5 en
escenarios adversariales. Un sistema puede aprender pistas superficiales que no
representan comprensión. En MAP se deben construir pruebas de contraste: añadir
palabras irrelevantes, corregir ortografía, invertir una negación y cambiar un
número. El modelo debería ser estable ante cambios inocuos y sensible ante cambios
matemáticamente decisivos.

% =============================================================================
\chapter{Síntesis comparativa de la evidencia}
% =============================================================================

\small
\begin{longtable}{@{}P{2.55cm}P{2.5cm}P{3.05cm}P{3.15cm}P{3.05cm}@{}}
\caption{Estudios de revista y aporte transferible al caso MAP}
\label{tab:sintesis}\\
\toprule
\textbf{Estudio} & \shortstack{\textbf{Problema/}\\\textbf{datos}} & \textbf{Método} &
\textbf{Hallazgo relevante} & \textbf{Implicación para MAP}\\
\midrule
\endfirsthead
\multicolumn{5}{l}{\sffamily\small\color{GrisTexto}Continuación de la Tabla~\ref{tab:sintesis}}\\
\toprule
\textbf{Estudio} & \shortstack{\textbf{Problema/}\\\textbf{datos}} & \textbf{Método} &
\textbf{Hallazgo relevante} & \textbf{Implicación para MAP}\\
\midrule
\endhead
\bottomrule
\endfoot

Burrows et al. (2015) & Revisión de más de 80 trabajos de respuestas cortas. &
Reglas, similitud, ML y evaluación. & Documenta la evolución del campo y la
escasez histórica de benchmarks públicos. & Definir baseline y evaluación
reproducible antes de aumentar complejidad.\\[2mm]

Sahu y Bhowmick (2020) & Calificación automática en conjuntos de referencia. &
Características léxicas, alineamiento, semántica y \emph{stacking}. & Las señales
complementarias mejoran a modelos aislados. & Ensamblar evidencia lineal y
contextual usando predicciones fuera de muestra.\\[2mm]

Zhang et al. (2020) & Respuestas estudiantiles condicionadas por pregunta y
estudiante. & Red de creencia profunda con modelos de contexto. & Pregunta y
estudiante agregan información frente a respuesta sola. & Concatenar pregunta,
opción y explicación; evaluar el aporte de cada bloque.\\[2mm]

Zhang et al. (2022) & Más de 16,000 respuestas semiabiertas, siete preguntas. &
LSTM con información general y específica del dominio. & El orden y el dominio
mejoran la representación. & Ablación secuencial opcional; preservar orden y
notación.\\[2mm]

Zhu et al. (2022) & SemEval 2013 y conjunto Mohler. & BERT, BiLSTM, cápsulas y
pérdida especializada. & El preentrenamiento contextual supera referencias en sus
conjuntos. & Ajustar un Transformer y compararlo bajo la misma partición.\\[2mm]

del Gobbo et al. (2023) & Múltiples conjuntos públicos de ASAG. & TF--IDF, BERT y
regresores. & La combinación léxica-semántica es robusta entre dominios. & Conservar
TF--IDF como baseline y componente de ensamble.\\[2mm]

Bonthu et al. (2023) & ASAG con datos etiquetados limitados. & Transferencia con
Transformers de oraciones y aumento. & Reporta +4.91\% en exactitud y +3.36\% en F1
en su escenario. & Probar aumento conservador, nunca alterar significado
matemático.\\[2mm]

Padó et al. (2024) & ASAG en alemán, preguntas vistas y no vistas. & Transformers y
modelos lingüísticos. & El Transformer no domina necesariamente en preguntas no
vistas. & La prueba agrupada por pregunta es indispensable.\\[2mm]

Bonthu et al. (2024) & Adaptación de dominio para respuestas cortas. &
Preentrenamiento específico y ajuste supervisado. & PyDSBERT reporta 88.77\% de
exactitud y 91.59\% F1 en SPRAG. & Probar preentrenamiento continuado con corpus
matemático.\\[2mm]

Schneider et al. (2023) & Cerca de 10 millones de pares, múltiples idiomas y tipos.
& Transformer con flujo humano. & La intervención humana controla fallos y eleva
confiabilidad. & Abstención por incertidumbre y revisión de casos críticos.\\[2mm]

Filighera et al. (2024) & Ataques de caja negra a ASAG. & Inserción adversarial de
adjetivos/adverbios contra BERT/T5. & Descensos de 10--22 puntos evidencian atajos
léxicos. & Incluir pruebas de contraste y auditoría de tokens influyentes.\\[2mm]

Grévisse (2024) & 2,288 respuestas, 82 preguntas, 12 cursos y tres idiomas. &
GPT--4 y Gemini con prompting. & Acuerdo moderado y necesidad de supervisión. & LLM
solo como comparador/auditor; registrar versión y prompts.\\[2mm]

Wang et al. (2025) & Clasificación de experiencia científica en respuestas
escritas. & BERT, ontología y balanceo. & El conocimiento estructurado complementa
la representación textual. & Usar descripciones y compatibilidades de la taxonomía
de errores.\\[2mm]

Otero et al. (2025) & 55 errores matemáticos y 220 ejemplos diagnósticos. & LLM con
restricción temática y revisión experta. & Hasta 83.9\% en la configuración
reportada. & Limitar candidatos por tema/pregunta y mantener revisión docente.\\[2mm]

Morley et al. (2026) & Revisión de 21 artículos sobre respuestas científicas. &
BERT a GPT--4. & Predominan BERT/variantes; contexto ayuda; persisten riesgos. &
Transformer contextual con métricas de confiabilidad y explicación.\\
\end{longtable}
\normalsize

\section{Patrones convergentes}

\begin{enumerate}
  \item \textbf{El contexto importa.} Pregunta, dominio y rúbrica contienen señal
  que la respuesta aislada no puede reconstruir (Zhang et al., 2020;
  Morley et al., 2026).
  \item \textbf{Lo híbrido es competitivo.} TF--IDF y BERT capturan fenómenos
  distintos; combinarlos mejora robustez y conserva interpretabilidad parcial
  (del Gobbo et al., 2023; Sahu \& Bhowmick, 2020).
  \item \textbf{El dominio merece adaptación.} Preentrenar o incorporar conocimiento
  específico ayuda cuando el vocabulario y las relaciones difieren del texto
  general (Bonthu et al., 2024; Wang et al., 2025).
  \item \textbf{La validación define la conclusión.} Una partición con preguntas
  vistas responde a una pregunta distinta de otra con preguntas no vistas (Padó
  et al., 2024).
  \item \textbf{Precisión no equivale a confiabilidad.} Los modelos pueden explotar
  atajos y fallar ante perturbaciones; la calibración, el análisis de error y la
  intervención humana son componentes del sistema (Filighera et al., 2024;
  Schneider et al., 2023).
\end{enumerate}

\section{Vacíos identificados}

La literatura de ASAG suele predecir una nota ordinal o binaria; MAP requiere
clasificar la naturaleza del razonamiento y ordenar pares jerárquicos. Muchos
estudios usan pocas preguntas, lo que dificulta estudiar transferencia. Además,
las comparaciones entre modelos rara vez controlan simultáneamente clases de cola
larga, notación matemática, calibración y robustez adversarial. Finalmente, los
LLM recientes añaden capacidad semántica, pero también dependencia de servicios,
versiones y políticas difíciles de congelar.

Estos vacíos hacen que la contribución del proyecto no sea simplemente ajustar un
modelo. El valor reside en evaluar un pipeline contextual bajo dos escenarios
complementarios ---preguntas conocidas y transferencia a preguntas no vistas---,
explicar la ganancia sobre una línea base transparente y documentar en qué
preguntas y conceptos falla.

% =============================================================================
\chapter{Selección preliminar de algoritmos}
% =============================================================================

\section{Criterios de decisión}

Se consideran siete criterios: adecuación a texto corto, uso del contexto,
desempeño esperado en cola larga, interpretabilidad, costo, facilidad de
reproducir y evidencia académica. La Tabla~\ref{tab:matriz-decision} resume un
juicio preliminar en escala 1--5. No son resultados del conjunto MAP; orientan qué
experimentos merecen realizarse.

\begin{table}[H]
\centering
\caption{Matriz cualitativa de decisión (1 = desfavorable; 5 = muy favorable;
Ctx. = contexto; Int. = interpretabilidad; Evid. = evidencia)}
\label{tab:matriz-decision}
\small
\begin{tabularx}{\textwidth}{@{}Y*{6}{C{1.25cm}}@{}}
\toprule
\textbf{Familia} & \textbf{Texto corto} & \textbf{Ctx.} &
\textbf{Cola larga} & \textbf{Int.} & \textbf{Costo} &
\textbf{Evid.}\\
\midrule
Reglas/similitud & 3 & 2 & 1 & 5 & 5 & 3\\
TF--IDF + logística & 4 & 3 & 3 & 5 & 5 & 5\\
TF--IDF + SVM lineal & 5 & 3 & 3 & 4 & 5 & 5\\
BiLSTM & 3 & 4 & 3 & 2 & 3 & 4\\
BERT/DeBERTa & 5 & 5 & 4 & 2 & 3 & 5\\
Modelo adaptado a matemáticas & 5 & 5 & 4 & 2 & 2 & 4\\
LLM por API & 4 & 5 & 3 & 2 & 1 & 3\\
Ensamble lineal + Transformer & 5 & 5 & 4 & 3 & 3 & 5\\
\bottomrule
\end{tabularx}
\end{table}

\section{Nivel 1: baseline lineal híbrido}

\nivel{Modelos seleccionados:} regresión logística multinomial y SVM lineal
One-vs-Rest sobre la unión de TF--IDF de palabras (1--2 gramas) y caracteres
(3--5 gramas).

\textbf{Razones para considerarlos.}

\begin{itemize}
  \item Son referencias sólidas en ASAG y permiten saber si un modelo complejo
  aprende más que vocabulario y frases frecuentes.
  \item Los caracteres toleran errores ortográficos y conservan formas parciales de
  números, fracciones y operadores.
  \item La regresión logística produce probabilidades; la SVM suele ser fuerte en
  espacios dispersos y puede calibrarse posteriormente.
  \item Los pesos de términos permiten inspeccionar atajos, palabras asociadas con
  cada concepto y posibles fugas por pregunta.
  \item Su costo bajo facilita validación cruzada, búsqueda de hiperparámetros y
  análisis por semillas.
\end{itemize}

El baseline debe comparar tres entradas: explicación sola; explicación + opción;
y pregunta + opción + explicación. Así se cuantifica cuánto aporta el contexto sin
atribuir automáticamente toda mejora al Transformer.

\section{Nivel 2: Transformer contextual y multitarea}

\nivel{Modelo principal:} DeBERTa-v3-base o, si el presupuesto computacional lo
exige, BERT/RoBERTa base. La entrada concatena pregunta, opción y explicación. La
salida emplea dos cabezas: categoría general y error conceptual condicional.

Sea $h_{\mathrm{CLS}}$ la representación contextual conjunta. Las dos
distribuciones se obtienen mediante

\[
p(c\mid x)=\operatorname{softmax}(W_c h_{\mathrm{CLS}}+b_c),\qquad
p(m\mid x)=\operatorname{softmax}(W_m h_{\mathrm{CLS}}+b_m).
\]

La pérdida multitarea puede escribirse como

\[
\mathcal{L}=\mathcal{L}_{cat}+\lambda\,I(m\text{ aplica})\mathcal{L}_{misc},
\]

donde la máscara $I(\cdot)$ evita convertir la ausencia estructural de
\codigo{Misconception} en una clase artificial. Se evaluarán pérdida ponderada,
\emph{focal loss} o muestreo balanceado únicamente con validación, pues una
corrección agresiva puede descalibrar el ranking.

\textbf{Razones para considerarlo.}

\begin{itemize}
  \item La literatura reciente muestra ventajas de representaciones
  bidireccionales preentrenadas en respuestas cortas (Zhu et al., 2022; Morley et
  al., 2026).
  \item La atención contextual puede distinguir negaciones y el papel de una frase
  dentro de la pregunta.
  \item La arquitectura multitarea refleja que categoría y error conceptual están
  relacionados, pero no son idénticos.
  \item DeBERTa es suficientemente establecido para clasificación y viable en una
  GPU académica con secuencias cortas.
\end{itemize}

\section{Nivel 3: adaptación al lenguaje matemático}

\nivel{Alternativas:}\par (a) inicializar desde MathBERT; o (b) continuar el
preentrenamiento de BERT/DeBERTa mediante \emph{masked language modeling} sobre
explicaciones y material matemático no etiquetado, y luego ajustar la tarea.

Esta fase se ejecutará después del Transformer base. Solo se justifica si aporta
una mejora estable en preguntas no vistas o clases raras. El corpus de adaptación
debe excluir cualquier texto de validación que contenga señales de etiqueta. Las
explicaciones sin etiqueta pueden participar en preentrenamiento auto-supervisado
si se respetan las particiones.

\textbf{Razones para considerarlo.} La adaptación de dominio ha mejorado ASAG en
literatura reciente (Bonthu et al., 2024), y el vocabulario matemático contiene
usos especializados. No se selecciona directamente como única solución porque
añade costo y existe riesgo de sobreajuste a las 15 preguntas disponibles.

\section{Nivel 4: ensamble calibrado y ranking}

La salida final exige parejas. Una aproximación coherente calcula

\[
p(c,m\mid x)=p(c\mid x)\,p(m\mid c,x),
\]

aplica una máscara de compatibilidad y ordena los tres pares con mayor
probabilidad. Para categorías sin concepto erróneo se utiliza el valor permitido
por el formato de la competencia. La mezcla de baseline y Transformer se ajusta
con predicciones fuera de muestra. Temperatura, Platt o regresión isotónica se
compararán para calibrar, sin usar el conjunto de prueba.

\textbf{Razones para considerarlo.} MAP@3 recompensa el orden y otorga crédito
cuando la etiqueta correcta aparece entre las primeras posiciones. Un ensamble
puede rescatar clases en las que el modelo lineal reconoce una expresión precisa y
otras en las que el Transformer entiende una paráfrasis. La calibración es clave
porque la combinación de probabilidades sin escala comparable es arbitraria.

\section{Modelo no seleccionado como solución principal}

Los LLM generativos por API se incluyen solo como referencia. Pueden formular una
explicación legible y usar descripciones de etiquetas, pero su salida puede variar
entre versiones y ejecuciones. También elevan costo y riesgo de privacidad y son
susceptibles a instrucciones contenidas en el texto estudiantil. Si se prueban,
se congelará un conjunto pequeño, se validará formato, se registrará el prompt y
se impedirá que texto del estudiante modifique instrucciones del sistema.

\begin{decision}[Selección final de la revisión]
\textbf{Baseline obligatorio:} TF--IDF de palabras + caracteres con logística y
SVM. \textbf{Modelo principal:} DeBERTa/BERT contextual multitarea.
\textbf{Experimento avanzado:} adaptación de dominio con MathBERT o preentrenamiento
continuado. \textbf{Salida operativa:} ensamble calibrado, restricciones
jerárquicas y ranking top--3. \textbf{LLM:} comparación diagnóstica, no producción.
\end{decision}

% =============================================================================
\chapter{Plan experimental propuesto}
% =============================================================================

\section{Particiones y control de fuga}

Se recomienda conservar un conjunto de prueba final sin contacto durante la
selección. Dentro del desarrollo se reportarán dos esquemas complementarios. La
validación cruzada estratificada por etiqueta compuesta medirá el desempeño sobre
preguntas conocidas; la validación agrupada por \codigo{QuestionId} estimará la
transferencia a preguntas no vistas. Como solo existen quince preguntas, se
usarán tres o cinco folds en el esquema agrupado y se documentará qué preguntas
quedan fuera en cada iteración. La selección operativa dará prioridad al primer
escenario, mientras que el segundo funcionará como prueba explícita de
extrapolación.

Los transformadores de texto, vocabularios, ponderaciones de clase, calibradores y
meta-modelos deben ajustarse solo con el fold de entrenamiento. Las predicciones
para el ensamble deben ser \emph{out-of-fold}. El texto duplicado o casi duplicado
debe permanecer en el mismo grupo para impedir contaminación.

\section{Métricas}

Si $r_i$ es la posición de la primera etiqueta correcta para el ejemplo $i$, la
precisión promedio hasta tres candidatos se expresa como

\[
\operatorname{MAP@3}=\frac{1}{N}\sum_{i=1}^{N}
\begin{cases}
1/r_i, & r_i\leq3,\\
0, & r_i>3\text{ o no aparece}.
\end{cases}
\]

MAP@3 será la métrica de selección por alinearse con el reto. Se acompañará de:

\begin{itemize}
  \item macro F1 y recall macro para visibilizar clases raras;
  \item F1 ponderado y exactitud como contexto, nunca de forma aislada;
  \item desempeño por \codigo{QuestionId}, categoría y concepto erróneo;
  \item recall@3 y rango recíproco medio;
  \item pérdida logarítmica y error de calibración esperado;
  \item media y desviación entre semillas y folds.
\end{itemize}

\section{Secuencia de experimentos}

\begin{table}[H]
\centering
\caption{Hoja de ruta experimental}
\label{tab:experimentos}
\begin{tabularx}{\textwidth}{@{}C{1.1cm}P{3.2cm}YY@{}}
\toprule
\textbf{Fase} & \textbf{Experimento} & \textbf{Pregunta que responde} &
\textbf{Criterio de avance}\\
\midrule
0 & Dummy por frecuencia y por pregunta & ¿Qué desempeño se obtiene sin leer el
texto? & Todo modelo real debe superarlo en MAP@3 y macro F1.\\
1 & TF--IDF palabra/carácter + logística/SVM & ¿Cuánta señal existe en expresiones
superficiales? & Baseline estable y errores interpretables.\\
2 & Ablación de entradas & ¿Qué aportan opción y pregunta? & Mejoras consistentes
en folds agrupados.\\
3 & BERT/DeBERTa contextual & ¿La semántica contextual supera al baseline? & Mejora
con costo y varianza aceptables.\\
4 & Pérdida multitarea y balanceo & ¿La jerarquía ayuda a conceptos raros? & Mayor
macro F1 sin degradar MAP@3/calibración.\\
5 & Adaptación matemática & ¿El dominio aporta transferencia? & Ganancia en
preguntas no vistas y clases raras.\\
6 & Ensamble y calibración & ¿Las señales son complementarias? & Mejora fuera de
muestra y probabilidades más confiables.\\
7 & Robustez y revisión humana & ¿El modelo usa contenido o atajos? & Estabilidad
ante perturbaciones inocuas y ruta de abstención.\\
\bottomrule
\end{tabularx}
\end{table}

\section{Hiperparámetros iniciales}

\begin{table}[H]
\centering
\caption{Espacio inicial de búsqueda, sujeto al presupuesto computacional}
\label{tab:hiperparametros}
\begin{tabularx}{\textwidth}{@{}P{3.6cm}YY@{}}
\toprule
\textbf{Componente} & \textbf{Valores iniciales} & \textbf{Observación}\\
\midrule
TF--IDF palabra & n-gramas (1,1), (1,2); min df 2--5; sublinear TF & No eliminar
negaciones; revisar \emph{stopwords}.\\
TF--IDF carácter & n-gramas (3,5), (3,6); límites de palabra y texto & Captura
ortografía y notación parcial.\\
Logística/SVM & $C\in\{0.1,0.5,1,2,5\}$; peso balanceado sí/no & Calibrar SVM en
fold separado.\\
Transformer & tasa $1$--$5\times10^{-5}$; 3--8 épocas; longitud 128--256 &
\emph{Early stopping} por MAP@3 agrupado.\\
Balanceo & ponderación inversa suavizada, focal $\gamma\in\{1,2\}$ & Comparar
calibración y clases frecuentes.\\
Multitarea & $\lambda\in\{0.25,0.5,1,2\}$ & Máscara para etiqueta no aplicable.\\
Ensamble & $\alpha\in\{0,0.1,\ldots,1\}$ & Ajustar solo con out-of-fold.\\
Semillas & mínimo 3 & Reportar promedio y dispersión.\\
\bottomrule
\end{tabularx}
\end{table}

\section{Análisis de errores y explicabilidad}

Se construirá una matriz de confusión para etiquetas con soporte suficiente y una
tabla de los pares confundidos. Para el modelo lineal se inspeccionarán n-gramas de
mayor peso. Para el Transformer se usarán técnicas de atribución como gradientes
integrados o eliminación de tokens, evitando interpretar pesos de atención como
explicación causal.

Cada muestra auditada incluirá pregunta, opción, explicación, etiqueta real, top--3,
probabilidades y observación del revisor. Se priorizarán: clases minoritarias,
predicciones confiadas e incorrectas, desacuerdos entre baseline y Transformer,
casos que cambian al modificar una negación y ejemplos sintéticos.

% =============================================================================
\chapter{Riesgos, ética y reproducibilidad}
% =============================================================================

\section{Sesgos y cobertura}

El conjunto MAP no representa automáticamente a todos los estudiantes, niveles,
idiomas o currículos. El desempeño puede variar con ortografía, dialecto,
competencia lingüística o formas distintas de explicar el mismo concepto. Una
etiqueta de ``concepto erróneo'' describe una respuesta, no una identidad estable
del estudiante. El sistema no debe usarse para decisiones disciplinarias ni para
inferir capacidad general.

\section{Incertidumbre y supervisión}

Se recomienda una política de abstención. Si la probabilidad está mal calibrada,
el margen entre los dos primeros candidatos es pequeño o la entrada queda fuera
de distribución, el sistema debe mostrar tres sugerencias y solicitar revisión.
La literatura de sistemas confiables favorece combinar automatización con control
humano, en lugar de presentar la predicción como verdad definitiva (Schneider et
al., 2023; Grévisse, 2024).

\section{Privacidad y seguridad}

Las explicaciones pueden contener datos personales introducidos accidentalmente.
Antes de usar servicios externos debe existir anonimización y autorización. Los
prompts no deben permitir que el contenido de la respuesta sustituya las
instrucciones de clasificación. Además, las pruebas adversariales propuestas por
Filighera et al. (2024) deben formar parte de la evaluación, no considerarse una
tarea posterior opcional.

\section{Reproducibilidad}

Todo experimento debe registrar versión del conjunto, hash de particiones,
semillas, librerías, hiperparámetros, hardware y tiempo. Se conservarán matrices de
resultados por fold y no solo el mejor número. El repositorio debe separar datos,
notebooks, código fuente, modelos y reportes; los archivos grandes o con licencia
restrictiva no deben subirse sin verificación.

% =============================================================================
\chapter{Conclusiones}
% =============================================================================

La revisión bibliográfica permite responder las dos actividades solicitadas. En
problemas similares se han empleado reglas y similitud, TF--IDF con regresión o
SVM, ensambles de características, LSTM, Transformers ajustados, adaptación de
dominio y LLM mediante prompting. No existe un ganador universal: el resultado
depende de si las preguntas son vistas, del tamaño del corpus, del dominio, de la
distribución de clases y del nivel de supervisión humana.

Para MAP, la evidencia favorece un enfoque contextual e híbrido. El baseline
TF--IDF proporciona rapidez, transparencia y control. El Transformer incorpora
pregunta, opción y explicación para capturar semántica. La formulación multitarea
respeta la estructura entre categoría y error conceptual. La adaptación de
dominio es un experimento avanzado justificado por el lenguaje matemático, y el
ensamble calibrado transforma las probabilidades en el ranking top--3 exigido.

La contribución más importante del diseño será una evaluación honesta. Comparar
particiones estratificada y agrupada, reportar macro F1 y MAP@3, analizar clases
raras, probar perturbaciones y documentar incertidumbre evitará que una puntuación
elevada por memorización se confunda con comprensión. Bajo estas condiciones, el
sistema puede apoyar al docente a localizar patrones de razonamiento y priorizar
retroalimentación, sin reemplazar su juicio.

\begin{decision}[Recomendación de implementación]
Iniciar con el baseline y las ablaciones de contexto; ajustar después un
DeBERTa/BERT multitarea; probar adaptación matemática solo si mejora de forma
consistente ambos escenarios de validación; y cerrar con calibración, ensamble y
auditoría humana. Esta secuencia
maximiza el aprendizaje experimental por unidad de costo y mantiene una referencia
interpretable en cada etapa.
\end{decision}

% =============================================================================
\backmatter
\chapter*{Referencias}
\addcontentsline{toc}{chapter}{Referencias}
\markboth{REFERENCIAS}{REFERENCIAS}

\small
\begin{list}{}{%
  \setlength{\leftmargin}{1.27cm}%
  \setlength{\itemindent}{-1.27cm}%
  \setlength{\itemsep}{0.9em}}

\item Bonthu, S., Sree, S. R., \& Krishna Prasad, M. H. M. (2023). Improving the
performance of automatic short answer grading using transfer learning and
augmentation. \emph{Engineering Applications of Artificial Intelligence, 123},
106292. \url{https://doi.org/10.1016/j.engappai.2023.106292}

\item Bonthu, S., Sree, S. R., \& Krishna Prasad, M. H. M. (2024). Framework for
automation of short answer grading based on domain-specific pre-training.
\emph{Engineering Applications of Artificial Intelligence, 137}, 109163.
\url{https://doi.org/10.1016/j.engappai.2024.109163}

\item Burrows, S., Gurevych, I., \& Stein, B. (2015). The eras and trends of
automatic short answer grading. \emph{International Journal of Artificial
Intelligence in Education, 25}(1), 60--117.
\url{https://doi.org/10.1007/s40593-014-0026-8}

\item del Gobbo, E., Guarino, A., Cafarelli, B., \& Grilli, L. (2023). GradeAid: A
framework for automatic short answers grading in educational contexts---Design,
implementation and evaluation. \emph{Knowledge and Information Systems, 65},
4295--4334. \url{https://doi.org/10.1007/s10115-023-01892-9}

\item Devlin, J., Chang, M.-W., Lee, K., \& Toutanova, K. (2019). BERT:
Pre-training of deep bidirectional transformers for language understanding. In
\emph{Proceedings of the 2019 Conference of the North American Chapter of the
Association for Computational Linguistics: Human Language Technologies} (Vol. 1,
pp. 4171--4186). Association for Computational Linguistics.
\url{https://doi.org/10.18653/v1/N19-1423}

\item Filighera, A., Ochs, S., Steuer, T., \& Tregel, T. (2024). Cheating automatic
short answer grading with the adversarial usage of adjectives and adverbs.
\emph{International Journal of Artificial Intelligence in Education, 34},
616--646. \url{https://doi.org/10.1007/s40593-023-00361-2}

\item Grévisse, C. (2024). LLM-based automatic short answer grading in
undergraduate medical education. \emph{BMC Medical Education, 24}, 1060.
\url{https://doi.org/10.1186/s12909-024-06026-5}

\item He, P., Liu, X., Gao, J., \& Chen, W. (2021). DeBERTa: Decoding-enhanced
BERT with disentangled attention. In \emph{International Conference on Learning
Representations}. \url{https://openreview.net/forum?id=XPZIaotutsD}

\item Kaggle. (2025). \emph{MAP---Charting student math misunderstandings}
[Data science competition].
\url{https://www.kaggle.com/competitions/map-charting-student-math-misunderstandings/overview}

\item Morley, F., Walland, E., \& Vidal Rodeiro, C. L. (2026). Auto-marking short
answer questions in science: The foundational years of transformer-based models
from BERT to GPT-4. \emph{International Journal of Artificial Intelligence in
Education, 36}(1--2), 100005.
\url{https://doi.org/10.1016/j.ijaied.2026.100005}

\item Otero, N., Druga, S., \& Lan, A. (2025). A benchmark for math
misconceptions: Bridging gaps in middle school algebra with AI-supported
instruction. \emph{Discover Education, 4}, 277.
\url{https://doi.org/10.1007/s44217-025-00742-w}

\item Padó, U., Eryilmaz, Y., \& Kirschner, L. (2024). Short-answer grading for
German: Addressing the challenges. \emph{International Journal of Artificial
Intelligence in Education, 34}(4), 1321--1352.
\url{https://doi.org/10.1007/s40593-023-00383-w}

\item Rittle-Johnson, B., Adler, R., Durkin, K., Burleigh, L., King, J., \&
Crossley, S. (2025). Detecting math misconceptions: An AI benchmark dataset. In
\emph{Proceedings of the 1st Workshop on Advancing Artificial Intelligence through
Theory of Mind} (pp. 20--24). Association for Computational Linguistics.
\url{https://aclanthology.org/2025.aimecon-wip.3/}

\item Rushton, S. J. (2018). Teaching and learning mathematics through error
analysis. \emph{Fields Mathematics Education Journal, 3}, Article 4.
\url{https://doi.org/10.1186/s40928-018-0009-y}

\item Sahu, A., \& Bhowmick, P. K. (2020). Feature engineering and ensemble-based
approach for improving automatic short-answer grading performance. \emph{IEEE
Transactions on Learning Technologies, 13}(1), 77--90.
\url{https://doi.org/10.1109/TLT.2019.2897997}

\item Schneider, J., Richner, R., \& Riser, M. (2023). Towards trustworthy
AutoGrading of short, multi-lingual, multi-type answers. \emph{International
Journal of Artificial Intelligence in Education, 33}(1), 88--118.
\url{https://doi.org/10.1007/s40593-022-00289-z}

\item Shen, J. T., Yamashita, M., Prihar, E., Heffernan, N. T., Wu, X., Lee, D., \&
Wang, D. (2021). MathBERT: A pre-trained language model for general NLP tasks in
mathematics education. \emph{arXiv}.
\url{https://doi.org/10.48550/arXiv.2106.07340}

\item Wang, H., Haudek, K. C., Manzanares, A. D., Romulo, C. L., Royse, E. A., \&
Azzarello, C. B. (2025). Extending a pretrained language model (BERT) using an
ontological perspective to classify students' scientific expertise level from
written responses. \emph{Disciplinary and Interdisciplinary Science Education
Research, 7}, 26. \url{https://doi.org/10.1186/s43031-025-00149-5}

\item Zhang, L., Huang, Y., Yang, X., Yu, S., \& Zhuang, F. (2022). An automatic
short-answer grading model for semi-open-ended questions. \emph{Interactive
Learning Environments, 30}(1), 177--190.
\url{https://doi.org/10.1080/10494820.2019.1648300}

\item Zhang, Y., Lin, C., \& Chi, M. (2020). Going deeper: Automatic short-answer
grading by combining student and question models. \emph{User Modeling and
User-Adapted Interaction, 30}(1), 51--80.
\url{https://doi.org/10.1007/s11257-019-09251-6}

\item Zhu, X., Wu, H., \& Zhang, L. (2022). Automatic short-answer grading via
BERT-based deep neural networks. \emph{IEEE Transactions on Learning Technologies,
15}(3), 364--375. \url{https://doi.org/10.1109/TLT.2022.3175537}

\end{list}
\normalsize

\chapter*{Apéndice A. Correspondencia entre evidencia y decisión}
\addcontentsline{toc}{chapter}{Apéndice A. Correspondencia entre evidencia y decisión}
\markboth{APÉNDICE A}{APÉNDICE A}

\begin{tabularx}{\textwidth}{@{}P{4.0cm}YY@{}}
\toprule
\textbf{Decisión} & \textbf{Evidencia principal} & \textbf{Comprobación local}\\
\midrule
Baseline TF--IDF/SVM & Burrows et al. (2015); del Gobbo et al. (2023) & Validación
doble, pesos de n-gramas y costo.\\
Entrada con pregunta y opción & Zhang et al. (2020); Morley et al. (2026) &
Ablación de tres configuraciones de entrada.\\
Transformer contextual & Zhu et al. (2022); Morley et al. (2026) & MAP@3 y macro F1
contra baseline, mismas particiones.\\
Adaptación de dominio & Bonthu et al. (2024); Wang et al. (2025) & Ganancia en
preguntas no vistas y clases raras.\\
Ensamble híbrido & Sahu y Bhowmick (2020); del Gobbo et al. (2023) & Mezcla con
predicciones out-of-fold y calibración.\\
Validación por pregunta & Padó et al. (2024) & GroupKFold por \codigo{QuestionId}.\\
Pruebas adversariales & Filighera et al. (2024) & Perturbaciones inocuas y cambios
semánticos mínimos.\\
Supervisión humana & Schneider et al. (2023); Grévisse (2024) & Umbral de abstención
y auditoría de alta incertidumbre.\\
\bottomrule
\end{tabularx}

\chapter*{Apéndice B. Lista de verificación antes de entrenar}
\addcontentsline{toc}{chapter}{Apéndice B. Lista de verificación antes de entrenar}
\markboth{APÉNDICE B}{APÉNDICE B}

\begin{enumerate}
  \item Congelar ambas particiones y verificar que ningún duplicado cruce entre
  entrenamiento y validación; en el esquema agrupado tampoco debe cruzar una
  pregunta.
  \item Confirmar la semántica exacta de categorías, valores nulos estructurales y
  formato de pares permitido.
  \item Construir dummy y baseline antes de ajustar Transformers.
  \item Conservar números, signos, operadores y negaciones en preprocesamiento.
  \item Ajustar vectorizador, ponderaciones y calibrador dentro de cada fold.
  \item Reportar MAP@3, macro F1, resultados por pregunta/clase y variación entre
  semillas.
  \item Auditar ejemplos confiados e incorrectos, clases raras y desacuerdos entre
  modelos.
  \item Documentar licencias, versiones, semillas, hardware, tiempos y huella de
  datos.
  \item Definir un umbral de abstención y la intervención humana esperada.
  \item No interpretar las puntuaciones de la matriz cualitativa como resultados
  empíricos.
\end{enumerate}

\vfill
\begin{center}
\color{GrisTexto}\sffamily\small
Documento preparado para el Proyecto 2 de CC3084 -- Data Science.\\
Universidad del Valle de Guatemala, 2026.
\end{center}

\end{document}
